%% sections/06-conclusion.tex

\section{Conclusion}
\label{sec:conclusion}

Multi-session TRPG campaign arcs fail to complete at high rates, and when they
complete, they often do so through DM improvisation rather than designed architecture.
The emotional payoffs that make long-form TRPG play valuable --- the arcs that land,
the finales that feel earned, the moments players cite years later --- are structurally
unreliable in the absence of explicit campaign design.

The campaign spine addresses this unreliability through five components that together
specify what the campaign is about (central question), how the party will come to
understand it (hint schedule), what constitutes completion at different engagement
levels (branching end-conditions), when each character's arc is foregrounded
(spotlight schedule), and what the finale's emotional content will come from
(PC-authored deliverable). Across 7 complete campaigns in the Marathon AI-simulated
TRPG workshop, these five components produced:

\begin{itemize}
  \item \textbf{7/7 campaign completions} (100\% completion rate; 49 sessions total)
  \item \textbf{5/7 Route D or E achievements} (71\%; the end-conditions requiring
    highest sustained party engagement)
  \item \textbf{28/28 hints delivered} (100\%; three campaigns used recovery paths)
  \item \textbf{7/7 PC-authored deliverables produced} verbatim at the finale
  \item \textbf{Mean sessions-to-first-binding-PASS: 4.5} (range: 4--6 across
    all campaign archetypes)
\end{itemize}

The spine is archetype-independent. It produced equivalent completion and
quality outcomes across grief-themed, covenant, faith/authority, siege-resource,
professional-code, duty-spine, and no-designed-stakes campaign types. The
strangers result (Campaign 7, Route E, no designed personal stakes) is the
strongest evidence for this claim: the spine's five components produced Route E
completion for a party with no grief-hooks, no professional codes, and no
institutional identities, through a lightweight trust-accumulation mechanism
calibrated to the archetype's emotional register.

The paper's primary contribution is not the individual components --- hint
scheduling, route branching, and PC spotlights appear in various forms in
practitioner TRPG design resources. The contribution is the demonstrated
\textit{system}: five components that interact specifically and predictably,
producing reliable emotional completeness across diverse campaign archetypes
when all five are implemented together. The interaction structure
(Table~\ref{tab:component-interaction}) is the design knowledge this paper
contributes: hints enable routes; spotlights load deliverables; end-conditions
determine hint delivery priority; central questions drive hint coherence.

We make the full campaign spine template available as supplementary material,
including the complete specifications for all 7 campaigns in the corpus. Our
hope is that campaign designers --- whether running AI-simulated or human-run
TRPG --- will find the five-component structure useful as a design constraint
that does not determine narrative content but provides the architectural
backbone without which narrative content cannot reliably converge into an
emotionally complete arc.
